% Section 6.4 — Main Classification Performance
% Presents the primary SDI-4 classification results for the proposed method.

\section{Main Classification Performance}
\label{sec:main-classification}

This section reports the primary classification results on the SDI-4 test partition. The proposed method (YOLO26m-cls + ASL-LDAM + SimAM-DCFR) is compared against the CE baseline (YOLO26m-cls with cross-entropy loss) under the same training protocol.

\subsection{Training Configuration}
\label{sec:training-config}

Both the proposed method and the CE baseline were trained on the SDI-4 training partition (804 images) with the following shared configuration:

\begin{itemize}
    \item Input resolution: 224 $\times$ 224 pixels.
    \item Epochs: 30, with early stopping patience of 15.
    \item Batch size: 32.
    \item Optimiser: AdamW, initial learning rate 0.00125, cosine annealing schedule, final learning rate factor 0.01.
    \item Mixed precision: AMP enabled.
    \item Data augmentation: RandAugment, random erasing with probability 0.4.
    \item Workers: 4.
    \item Seed: 42.
\end{itemize}

The proposed method additionally applies:

\begin{itemize}
    \item \textbf{ASL-LDAM}: Asymmetric loss with label-distribution-aware margin. Configuration: $\gamma_{\text{pos}} = 0$, $\gamma_{\text{neg}} = 4$, label smoothing 0.1, LDAM max margin 0.5, LDAM scale 30. Class counts: $[282, 139, 229, 154]$.
    \item \textbf{SimAM-DCFR}: Simple attention module with dual-channel feature recalibration, inserted immediately before the original classification head. The attention branch is parameter-free; the texture branch and channel gate are trainable. Fusion: $\mathbf{F}' = \mathbf{F} + 0.5 \cdot \mathbf{G} \odot (\mathbf{F}_{\text{SimAM}} + \mathbf{F}_{\text{texture}})$, where $\mathbf{G}$ is a learned channel-gating vector. The SimAM energy parameter is $\lambda = 10^{-4}$.
\end{itemize}

\subsection{Primary Results}
\label{sec:primary-results}

The primary SDI-4 test results are reported in Table~\ref{tab:primary-results}.

\begin{table}[htbp]
\centering
\caption{Primary SDI-4 classification results (seed 42).}
\label{tab:primary-results}
\begin{tabular}{l C{2.2cm} C{2.2cm} C{2.2cm}}
\toprule
\textbf{Method} & \textbf{Macro-F1 (\%)} & \textbf{Accuracy (\%)} & \textbf{Kappa} \\
\midrule
CE Baseline (YOLO26m-cls)      & 89.02 & 89.02 & 0.8505 \\
Proposed (ASL-LDAM + SimAM-DCFR) & 91.01 & 91.33 & 0.8816 \\
\midrule
$\Delta$ & +1.99 & +2.31 & +0.0311 \\
\bottomrule
\end{tabular}
\end{table}

The proposed method improves upon the CE baseline by 1.99 percentage points in Macro-F1, 2.31 percentage points in accuracy, and 0.0311 in Cohen's Kappa. These improvements are observed under the recorded protocol and seed 42; they support the effectiveness of the ASL-LDAM loss and SimAM-DCFR attention module for this four-class shrimp disease classification task.

\subsection{Per-Class Analysis}
\label{sec:per-class}

The per-class F1 scores for both methods are available from the retained confusion matrices and classification reports in the experiment archives. The SDI-4 test partition comprises 61 Healthy, 29 BG, 50 WSSV, and 33 WSSV\_BG images. The primary sources of confusion are expected between WSSV and WSSV\_BG (co-infection boundary) and between Healthy and disease classes (false-positive/negative asymmetry).

\subsection{Model Capacity and Footprint}
\label{sec:model-capacity}

The proposed method introduces 529,408 trainable parameters through the SimAM-DCFR module, increasing the total parameter count from 10.35M (YOLO26m-cls backbone) to approximately 10.88M. The FP32 model size of the proposed classifier is approximately 39.52 MiB in TFLite format, as discussed in Section~\ref{sec:mobile}.

The SDI-4 trained checkpoint hash for the proposed method is:
\path{bf4d969465025e870a61be5e218dfc55b9b2e962dec416f220b1f130703fb1c7}

The CE baseline checkpoint hash is:
\path{4305e49158129c4a6acaa8fdaf7b5982a7f4d93cc233f11d17479e04b0e438c1}